%add all your local new commands to this file

\newcommand{\smiley}{:)}

\renewbibmacro*{index:name}[5]{%
  \usebibmacro{index:entry}{#1}
    {\iffieldundef{usera}{}{\thefield{usera}\actualoperator}\mkbibindexname{#2}{#3}{#4}{#5}}}

% \newcommand{\noop}[1]{}

\newcommand{\zB}{z.\,B.\ }

\newcommand{\wzB}{w.\,z.\,B.\ }

\newcommand{\ua}{u.\,a.\ }

\newcommand{\fbmarpar}[1]{\marginpar{\footnotesize{#1}}}


%\newcommand{\dh}{d.\,h.\ }

% http://tex.stackexchange.com/questions/230300/doing-something-like-psframebox-in-tikz#230306
\tikzset{
frbox/.style={
  rounded corners,
  draw,
  thick,
  inner sep=5pt
  }
}
\newcommand\TZbox[1]{\tikz\node[frbox,baseline] {#1};}

\newcommand{\gqq}[1]{\glqq{}#1\grqq{}}		%double
\newcommand{\gq}[1]{\glq{}#1\grq{}}			%simple

\newcommand{\blau}[1]{\textcolor{blue}{#1}}

%von Schäfer
\forestset{
  narroof/.style={roof, inner xsep=-0.25em, rounded corners},
}

% \newcommand{\tblsthickline}{{\color{gray}\rule{\textwidth}{1.5mm}}}

\newcounter{dgvertief} 
\newenvironment{Sandwich}[1]
  {
    \par\vspace{5mm}\noindent\tblsthickline
    {\par\vspace{3mm}\noindent\sffamily\large\bfseries#1\vspace{4mm}}
  }
  {\par\vspace{3mm}\noindent\tblsthickline\par\vspace{5mm}}


\DeclareNewTOC[
  type=exkurs,
  types=exkurse,
  name=Exkurs,
  listname={Liste der Exkurse}
]{loe}


\newcounter{dgexkurs}
\newenvironment{Exkurs}[1]
  {
    \refstepcounter{dgexkurs}
    %\begin{Sandwich}{#1\hfill Exkurs~\thechapter.\arabic{dgexkurs}}
    \begin{Sandwich}{#1\hfill Exkurs~\arabic{dgexkurs}}
    \addcontentsline{loe}{exkurs}{\protect\numberline{\thedgexkurs}#1}%
  }
{\end{Sandwich}}

% \newcounter{dgdef}
\newcounter{dgdef}[chapter]
\renewcommand{\thedgdef}{\thechapter.\arabic{dgdef}}
\newenvironment{Definition}[1]
  {
    \refstepcounter{dgdef}
    \begin{Sandwich}{#1\hfill Definition~\thechapter.\arabic{dgdef}}
  }
    {\end{Sandwich}}

\newcounter{dgLoesung}
\newenvironment{Loesung}[1]
  {
    \refstepcounter{dgLoesung}
    \begin{Sandwich}{#1}

    %\begin{Sandwich}{#1\hfill Loesung~\thechapter.\arabic{dgLoesung}}
    %\begin{Sandwich}{#1\hfill Loesung~\thechaptr}
    
  }
{\end{Sandwich}}

 % {\end{Sandwich}}

%für Silbengelenk, von Schäfer
\newcommand{\SonDiag}[2][0]{%
  \begin{tikzpicture}
    \textsave{.}
    \tikzset{
      normalseg/.style={fill=white},
      extrasyll/.style={circle, draw, fill=white},
      sylljoint/.style={diamond, draw, fill=white}
    }
    \node at (0,\plo) {P};
    \node at (0,\fri) {F};
    \node at (0,\nas) {N};
    \node at (0,\liq) {L};
    \node at (0,\vok) {V};

    % Draw the helper lines if required.
    \ifthenelse{\equal{#1}{0}}{}{%
      \foreach \y in {\plo, \fri, \nas, \liq,\vok} {%
	\draw [dotted, |-|] (0.25, \y) -- (#1.75, \y);
      }
    }

    \foreach [count=\x from 1, remember=\x as \lastx] \p / \y / \g in #2 {
      \ifthenelse{\equal{\y}{-1}}{\textsave{.}}{%

	% Draw the node, either plain, as Silbenbgelenk, or as extrasyllabic.
        \ifthenelse{\equal{\g}{1}}{%
	  \node (\x) [sylljoint] at (\x, \y) {\p};
	}{%
	  \ifthenelse{\equal{\g}{2}}{%
	    \node (\x) [extrasyll] at (\x, \y) {\p};
	  }{%
	    \node (\x) [normalseg] at (\x, \y) {\p};
	  }
	}

	% Draw the connection unless the previous node was not or was empty.
	\ifthenelse{\NOT\equal{\lastsaved}{.}}{%
	  \draw [->] (\lastx) to (\x);
	}{}
	\textsave{1}
      }
    }
    \aeundefinenumericnodes
  \end{tikzpicture}
}

\forestset{
  Ephr/.style={draw, ellipse, thick, inner sep=2pt},
  Eobl/.style={draw, rounded corners, inner sep=5pt},
  Eopt/.style={draw, rounded corners, densely dashed, inner sep=5pt},
  Erec/.style={draw, rounded corners, double, inner sep=5pt},
  Eoptrec/.style={draw, rounded corners, densely dashed, double, inner sep=5pt},
  Ehd/.style={rounded corners, fill=gray, inner sep=5pt,
    delay={content=\whyte{##1}}
  },
  Emult/.style={for children={no edge}, for tree={l sep=0pt}},
  phrasenschema/.style={for tree={l sep=2em, s sep=2em}},
  decide/.style={draw, chamfered rectangle, inner sep=2pt},
  finall/.style={rounded corners, fill=gray, text=white},
  intrme/.style={draw, rounded corners},
  yes/.style={edge label={node[near end, above, sloped, font=\scriptsize]{Ja}}},
  no/.style={edge label={node[near end, above, sloped, font=\scriptsize]{Nein}}},
  sake/.style={tier=preterminal},
  ake/.style={
    tier=preterminal
    },
}


%fuer Silbenboegen
% usage: \uarc[sep][depth]{text}
\NewDocumentCommand\uarc{O{.2ex} O{.3ex} m}{%
    \tikz[baseline=(arced node.base)]{
        \node [inner sep=0pt, outer sep=0pt] (arced node) {#3};
        \draw [transform canvas={yshift=-#1}] (arced node.south west) parabola [parabola height=-#2] (arced node.south east);
    }%
}


